\documentclass[11pt,a4paper]{article}
\usepackage[margin=25mm]{geometry}
\usepackage{amsmath,amssymb,array,hyperref}
\hypersetup{colorlinks=true,urlcolor=blue,linkcolor=blue}
\title{Literature and publication assessment\\Length-scaled rigid-body contact notation}
\author{Physics-engine project review}
\date{6 October 2026}
\begin{document}
\maketitle

\section{Assessment}
The length-scaling operation is not unique. A direct planar velocity/force precedent exists in 2011; a related spatial deformation/force precedent exists in 2014. The exact symbol arrangement $(\mathbf v,\ell\boldsymbol\omega)$ and $(\mathbf p,\mathbf L/\ell)$ was not located as a pair in this search. That absence is not evidence of priority: the momentum expression follows from established dual scaling. Translation-first ordering and additive contact matrices provide a clear presentation, but do not establish new mechanics.

The current article can be circulated as an explanatory technical note or tutorial with reproducible calculations and proper attribution. A research-paper contribution has not yet been established. Potential research directions are identified below; they are recommendations, not completed results or guarantees of acceptance.

\section{Closest notation precedents}
Vose, Umbanhowar and Lynch's RSS 2011 paper defines, in its appendix footnote 1 on PDF page 7, a planar twist whose angular-velocity component is multiplied by the radius of gyration $\rho$, and a wrench whose torque component is divided by $\rho$\cite{vose}. In our component names this reads
\begin{equation}
\boldsymbol\xi=\begin{bmatrix}v_x\\v_y\\\rho\omega\end{bmatrix},\qquad
\mathbf w=\begin{bmatrix}f_x\\f_y\\\tau/\rho\end{bmatrix}.
\end{equation}
It uses the same planar scaling operation as our notation, with a particular geometric/inertial choice of length. Its force vector is not literally our momentum vector.

Zhang, Su and Liao (2014), section 3.3, equation 21 of the author manuscript, multiply rotational deformation by a characteristic length and divide applied moment by that length\cite{zhang}. Their motion vector is rotation first and concerns small deformation. This is substantial prior art for the dual length scaling, rather than a claim that their entire collision formulation is ours.

For fixed $\ell$, integrating the scaled wrench through an impulse gives
\begin{equation}
\int\begin{bmatrix}\mathbf f\\\boldsymbol\tau/\ell\end{bmatrix}\,dt
=\begin{bmatrix}\Delta\mathbf p\\\Delta\mathbf L/\ell\end{bmatrix}.
\end{equation}
Thus the momentum pairing is a direct dual counterpart. Featherstone's spatial-vector documentation already treats motion vectors separately from their dual force, impulse and momentum vectors\cite{featherstone}. This inference does not depend on finding the exact letters $\ell$ and $\mathbf L$ in an earlier paper.

\clearpage
\section{Which parts are established?}
\begin{center}\small
\begin{tabular}{p{45mm}p{102mm}}
\hline
Our feature & Literature finding and implication\\
\hline
Length-scaled velocity and dual momentum & Direct velocity/wrench scaling in Vose et al.; related spatial scaling in Zhang et al.; momentum follows by duality. Scaling itself is not a new research result.\\[3pt]
Plus-sign contact map & $\boldsymbol\omega\times\mathbf r=(-\mathbf r)\times\boldsymbol\omega$ is cross-product antisymmetry. $\mathbf C=[\mathbf1\ {-\,[\mathbf r]_\times/\ell}]$ is a scaled contact Jacobian. Moving the sign into a named block does not change the algebra.\\[3pt]
Normal and tangential restitution plus friction & Already present in earlier impact models. Chatterjee and Ruina (1998) explicitly propose a 3D point-impact law with two restitution parameters and admissible frictional impulses\cite{chatterjee}. Their impulse-selection law is not asserted to be identical to ours.\\[3pt]
Coupled contact effective mass & $\mathbf W=\mathbf J\mathbf M^{-1}\mathbf J^T$ is the established Delassus operator. Acary and Collins-Craft discuss this and normal/tangential restitution in their 2025 paper, section 2.5\cite{acary}.\\[3pt]
Matrix-free contact evaluation & The SIGGRAPH 2022 course explicitly teaches coupled matrix-free evaluation using contact Jacobians and body inverse masses\cite{siggraph}. Our scatter/solve/gather organization is not sufficient evidence of a new algorithm.\\[3pt]
Physical energy acceptance & Dissipation problems with restitution and coupled contacts have prior analysis\cite{acary}. Checking actual kinetic change is useful implementation practice; the check alone is not a new constitutive law or general existence theorem.\\[3pt]
Experimental comparison and parameter interpretation & Fazeli et al. (2017) compare six planar impact models and examine fitted-parameter interpretation\cite{fazeli}. A comparison study needs a distinct question and evidence beyond repeating that general task.\\
\hline
\end{tabular}\end{center}

The combined 2D/3D derivation, notation choices, implementation and comparison figures are our authored work. Original authorship of an implementation or explanation is distinct from scientific novelty. The two-coordinate deviation map is useful communication, but does not retain the full signed motion state and is not established here as a novel error metric.

\section{What has not been proved by this review?}
This is a targeted public-literature review, not an exhaustive bibliography or priority certification. A search hit is not proof of identical models: coefficient definitions, contact regimes, friction laws, energy assumptions and simultaneous-impact rules must be compared separately. Nor does this review prove that every production-solver detail has appeared before. It establishes strong prior art for the article's principal mathematical ingredients; algorithm-level novelty remains unverified.

\clearpage
\section{What could support a research paper?}
\textbf{The strongest direction for the user's stated aim is independent prediction from documented material-pair inputs.} Pose a falsifiable question: with restitution and friction fixed from separate characterization, how accurately can the same model predict outgoing translation and spin over different initial states, sizes and non-spherical contact geometries?

Before claiming that contribution, establish all of the following:
\begin{enumerate}
\item Parameter characterization and validation impacts are independent. No coefficients are optimized against validation outcomes, and no unknown coefficients are silently supplied.
\item Shape, mass, full inertia, contact geometry, signed incoming/outgoing motion and measurement uncertainty are available at the level required by the model. Missing information restricts the claim.
\item Compare appropriate established impact laws on the same impacts, with parameter definitions and calibration rules declared. Publish component errors, spin/velocity direction errors where measured, energy balance and worst cases, not only the compact plot.
\item Preserve failure cases, uncertainty and validity ranges. Determine whether disagreement is due to missing geometry, measurement uncertainty, coefficient transfer, solver error or inadequacy of the contact law.
\end{enumerate}
Such a study could contribute a reproducible benchmark, a demonstrated transfer result, or a quantitative limit on what public material data can predict. No particular positive outcome is assumed.

An alternative is a specific solver or constitutive-law contribution: for example, a well-defined simultaneous-contact method with a new justified admissibility/selection rule, appropriate existence or convergence analysis, and successful comparisons. Rejection of energy-producing candidates is not the same as constructing a solution for all admissible inputs. Total-system dissipation also does not by itself establish local thermodynamic admissibility. Numerical benchmarks must establish accuracy and robustness; the user has removed the requested $2\times$ performance gate.

\section{Current readiness}
\begin{center}\small
\begin{tabular}{p{43mm}p{104mm}}
\hline
Existing evidence & What it supports now\\
\hline
24 Cornell glass impacts & Fixed published-input reproduction; normal RMSE $0.03667\,{\rm m/s}$, COM tangential RMSE $0.01780\,{\rm m/s}$. Characterization and worksheet may overlap; spin is reconstructed. Not a blind full-state validation.\\[3pt]
81 material-profile controls & Native analytic/energy integration checks. Synthetic controls do not establish material transfer to measured reality.\\[3pt]
Worked 3D cuboid & Reproducible coupled matrix calculation and coordinate-invariance checks. Illustrative prescribed inputs, not measured impact validation.\\[3pt]
Rubber and rocks & Conditional ball comparisons and historical approximate-shape rock fits. Complete independently documented same-pair inputs and full rock geometry/state remain missing.\\[3pt]
Larger coupled scenes & The global working/trajectory-accuracy baseline remains unqualified. No general robustness or speed advantage is claimed.\\
\hline
\end{tabular}\end{center}
\textbf{Recommendation:} publish the derivation as an attributed explanatory resource if desired; develop the research submission around one independently specified, held-out material/shape comparison first. The current evidence does not justify a claim of a new general collision theory or universally authentic material lookup.

\clearpage
\section{Search and source record}
Queries covered scaled twists/wrenches, characteristic length, angular momentum and homogeneous inertia, normal/tangential restitution, coupled impact energy, matrix-free Delassus solvers and empirical impact-model comparisons. References were examined through publisher pages, proceedings, institutional author material and author-uploaded manuscripts. Source PDF hashes and exact locations are retained in the accompanying \texttt{review.json}; copyrighted source PDFs are kept outside the repository. New search findings have also been incorporated into the derivation article's prior-art discussion.

\subsection*{Recent related benchmark}
The August 2026 GAUGE preprint supplies measured trajectories and calibrated physical metadata for a broader simulator benchmark\cite{gauge}. Appendix A.2.1 describes friction measurement and rail-collision restitution characterization. A general measured-physics benchmark is therefore not an unexplored publication category. We have not verified a separate signed tangential restitution profile or full spin-state suitability in its released data, and have not run our model on GAUGE. It is a candidate source and comparison target, not completed validation.

\begin{thebibliography}{9}
\bibitem{vose} T. H. Vose, P. Umbanhowar and K. M. Lynch. \emph{Sliding Manipulation of Rigid Bodies on a Controlled 6-DoF Plate}. Robotics: Science and Systems VII, 2011. Appendix footnote 1, PDF p7. \url{https://www.roboticsproceedings.org/rss07/p43.pdf}. DOI: \url{https://doi.org/10.15607/RSS.2011.VII.043}.
\bibitem{zhang} Y. Zhang, H.-J. Su and Q. Liao. \emph{Mobility criteria of compliant mechanisms based on decomposition of compliance matrices}. Mechanism and Machine Theory 79 (2014), 80--93. Author manuscript section 3.3, equation 21. \url{https://doi.org/10.1016/j.mechmachtheory.2014.04.010}.
\bibitem{featherstone} R. Featherstone. \emph{General Comments on Coordinate Transforms}. Spatial Vector and Rigid-Body Dynamics documentation. \url{https://royfeatherstone.org/spatial/v2/xforms.html}.
\bibitem{chatterjee} A. Chatterjee and A. Ruina. \emph{A New Algebraic Rigid-Body Collision Law Based on Impulse Space Considerations}. Journal of Applied Mechanics 65 (1998), 939--951. \url{https://doi.org/10.1115/1.2791938}.
\bibitem{acary} V. Acary and N. A. Collins-Craft. \emph{On the Moreau--Jean scheme with the Fr\'emond impact law: energy conservation and dissipation properties for elastodynamics with contact, impact and friction}. JTCAM, 2025. Section 2.5, equations 44--46. \url{https://jtcam.episciences.org/15960/pdf}. DOI: \url{https://doi.org/10.46298/jtcam.13480}.
\bibitem{siggraph} \emph{Contact and Friction Simulation for Computer Graphics}. SIGGRAPH 2022 course notes. Figure 43 and matrix-free contact tutorial, PDF pp163--164. \url{https://siggraphcontact.github.io/assets/files/SIGGRAPH22_friction_contact_notes.pdf}.
\bibitem{fazeli} N. Fazeli, S. Zapolsky, E. Drumwright and A. Rodriguez. \emph{Fundamental Limitations in Performance and Interpretability of Common Planar Rigid-Body Contact Models}. 2017, arXiv:1710.04979. \url{https://arxiv.org/abs/1710.04979}.
\bibitem{gauge} S. Wang et al. \emph{GAUGE: A Measurement-Grounded Benchmark for Physical Fidelity in Simulation Engines and Video World Models}. Preprint submitted 6 August 2026. Appendix A.2.1. \url{https://arxiv.org/abs/2608.05948}; \url{https://arxiv.org/html/2608.05948v1}.
\end{thebibliography}
\end{document}
